\chapter{Despliegue y uso del sistema}
\label{ap:despliegue}

El capítulo~\ref{cap:desarrollo} explica \emph{por qué} el despliegue tiene la forma que
tiene. Este apéndice explica \emph{cómo} se ejecuta, paso a paso, y está escrito para quien
tenga que ponerlo en marcha en un centro sin haber participado en el desarrollo: el
responsable informático el primer día, y el docente cada mañana.

\section{Qué hace falta antes de empezar}

El sistema se despliega en \textbf{un solo equipo del centro}, que actúa de servidor y hace
todo el trabajo. Los dispositivos del alumnado no necesitan absolutamente nada salvo un
navegador y estar en la misma red: no instalan, no calculan y no conceden permisos. Esa
asimetría es deliberada y se justifica en la sección~\ref{sec:arquitectura-app}; para el
despliegue lo que implica es que solo hay un equipo del que preocuparse.

Los requisitos de ese equipo no son una lista de recomendaciones: cada uno procede de una
medición del capítulo~\ref{cap:desarrollo} y puede comprobarse antes de dar una clase con el
sistema. Se presentan en dos tablas porque responden a dos preguntas distintas: la
tabla~\ref{tab:requisitos-despliegue} dice \emph{qué equipo hace falta}, que es una decisión
de compra, y la tabla~\ref{tab:requisitos-entorno} \emph{qué hay que preparar} en el que ya
se tenga.

\begin{table}[!ht]
\centering
\footnotesize
\setlength{\tabcolsep}{4pt}
\begin{tabular}{|p{2.5cm}|p{3.3cm}|p{6.2cm}|}
\hline
\textbf{Elemento} & \textbf{Requisito} & \textbf{De dónde sale} \\
\hline
Tarjeta gráfica &
    Dedicada, con al menos \textbf{4~GB} de memoria propia. No hace falta gama alta. &
    Medido: el modelo ocupa 1,60~GB de memoria de vídeo en el peor caso
    (tabla~\ref{tab:huella}). El resto es margen para el escritorio del propio equipo y la
    fragmentación de memoria. Sin tarjeta dedicada el sistema arranca, pero no sostiene el
    ritmo del audio en directo (sección~\ref{sec:hardware}). \\
\hline
Memoria del sistema &
    \textbf{8~GB}. &
    El proceso de reconocimiento ocupa 2,76~GB residentes; la aplicación y el sistema
    operativo caben holgadamente en el resto. \\
\hline
Procesador &
    Cuatro núcleos o más. Cualquiera de los últimos años sirve. &
    No calcula el modelo: despacha el trabajo a la tarjeta gráfica. Conviene saber que
    \textbf{un núcleo queda ocupado al 100\% de forma permanente} mientras el servicio está
    vivo, incluso en reposo: es espera activa del entorno de ejecución de la GPU, no trabajo
    real, y no indica ningún problema. \\
\hline
Almacenamiento &
    Depende de la vía: \textbf{5~GB} nativa, \textbf{100~GB} en contenedores sobre AMD. &
    El modelo son 2,85~GB en ambos casos. La diferencia no la causa el sistema sino la
    imagen base de la plataforma de AMD, y se detalla más abajo porque es lo bastante grande
    como para condicionar la elección de equipo. \\
\hline
\end{tabular}
\caption{Requisitos de hardware del equipo servidor, con la medición que respalda cada uno.}
\label{tab:requisitos-despliegue}
\end{table}

\begin{table}[!ht]
\centering
\footnotesize
\setlength{\tabcolsep}{4pt}
\begin{tabular}{|p{2.5cm}|p{3.3cm}|p{6.2cm}|}
\hline
\textbf{Elemento} & \textbf{Requisito} & \textbf{Observación} \\
\hline
Sistema operativo &
    Linux permite el despliegue en contenedores; en macOS con Apple Silicon hay que
    ejecutar de forma nativa. &
    No es indiferente: en macOS los contenedores no alcanzan la GPU. Conviene leer la
    sección~\ref{sec:contenedores} \emph{antes} de elegir procedimiento. \\
\hline
Acceso a la GPU &
    Con AMD, pertenecer a los grupos \texttt{video} y \texttt{render}. Con NVIDIA, el
    conjunto de herramientas de contenedores del fabricante. &
    Es el fallo de arranque más frecuente. Se comprueba con
    \texttt{getent group video render}, cuyos identificadores numéricos deben coincidir con
    los del descriptor de contenedores. \\
\hline
Micrófono &
    De solapa o diadema, próximo a la boca del docente. No vale el micrófono integrado del
    equipo ni uno ambiental. &
    Medido: el sistema tolera ruido hasta 15~dB de relación señal-ruido, y por debajo se
    degrada de golpe en lugar de gradualmente (sección~\ref{sec:ruido}). La proximidad del
    micrófono es lo que mantiene ese margen cuando el aula se anima. \\
\hline
Red &
    Red local del centro, con el equipo servidor y los dispositivos del alumnado en el
    mismo rango. &
    La aplicación rechaza toda conexión que no provenga de un rango privado. No es
    configurable: es un requisito legal del sistema
    (sección~\ref{sec:requisitos-legales}). Si hiciera falta acceso desde fuera, la vía es
    una red privada virtual hasta el centro. \\
\hline
Navegador &
    Cualquiera moderno, tanto en el equipo del docente como en los del alumnado. &
    La emisión debe abrirse en el propio equipo servidor, por una restricción de los
    navegadores que se explica en el paso 4. \\
\hline
\end{tabular}
\caption{Requisitos de entorno y periféricos.}
\label{tab:requisitos-entorno}
\end{table}

Conviene señalar lo que \emph{no} aparece en ninguna de las dos, porque suele darse por supuesto: no
hace falta conexión a internet durante la clase, ni cuenta en ningún servicio, ni licencia de
software. La única descarga es la del modelo, y ocurre una sola vez en el primer arranque.
Esto es consecuencia directa de que el reconocimiento se ejecute en el centro y no en la nube
(sección~\ref{sec:requisitos-legales}).

\subsection*{De dónde salen esas cifras}

Los tres requisitos de memoria y almacenamiento proceden de una medición directa del
servicio en ejecución, no del tamaño nominal del modelo, que es lo que suele citarse y
además induce a error. La tabla~\ref{tab:huella} recoge el detalle.

% GENERADO AUTOMATICAMENTE por research/eval/report/tabla_huella.py
% No editar a mano.
\begin{table}[H]
\centering
\small
\begin{tabular}{|l|r|r|r|}
    \hline
    \textbf{Modelo} & \textbf{Disco (GB)} & \textbf{VRAM pico (GB)} &
    \textbf{RAM del proceso (GB)} \\
    \hline
    \texttt{whisper-small} & 0.90 & 0.56 & 2.57 \\
    \texttt{whisper-medium} & 2.85 & 1.60 & 2.76 \\
    \texttt{whisper-medium + LoRA} & 2.85 & 1.63 & 2.78 \\
    \hline
\end{tabular}
\caption{Huella de recursos del servicio de reconocimiento, medida sobre
AMD Radeon RX 9070 XT con precisión \texttt{float16} y una ventana completa de 30~s de audio, que
es el caso peor. El pico de VRAM mayor observado es 1.63~GB, de donde sale el
requisito de 4~GB de la tabla~\ref{tab:requisitos-despliegue}: el margen cubre el
gestor de ventanas del propio equipo y la fragmentación de memoria.}
\label{tab:huella}
\end{table}


Dos lecturas de esa tabla importan para decidir la compra. La primera es que el requisito de
memoria de vídeo es \textbf{mucho menor de lo que sugiere el tamaño del modelo en disco}:
2,85~GB ocupados en almacenamiento se traducen en 1,60~GB de memoria de vídeo, porque los
pesos se cargan en precisión reducida. Una tarjeta de 4~GB, que hoy es la gama de entrada,
sobra. Lo que el sistema necesita no es una tarjeta grande sino, simplemente, una tarjeta.

La segunda es que existe una salida si el equipamiento disponible se queda corto: el modelo
inmediatamente menor baja el pico a 0,56~GB y el almacenamiento a 0,90~GB. Cuesta calidad,
entre uno y dos puntos de tasa de error según el material
(tabla~\ref{tab:baseline-modelos-tedx-es}), pero es una degradación conocida y medida, no una
sorpresa. La diferencia entre ambos modelos se estrecha además precisamente en habla
espontánea, que es el registro del aula.

Incorporar el modelo adaptado no cambia el cuadro: el adaptador añade 19~MB en disco y
30~MB de memoria de vídeo, porque no sustituye al modelo sino que le agrega unas matrices
pequeñas. Es la ventaja práctica de la técnica que ganó la comparativa, y la razón de que
pueda distribuirse por separado.

\subsection*{El almacenamiento, que es la sorpresa desagradable}

Las cifras anteriores se refieren al modelo. El despliegue completo en contenedores sobre
AMD ocupa mucho más, y conviene decirlo con claridad porque es el requisito que más se
subestima:

\begin{table}[H]
\centering
\footnotesize
\begin{tabular}{|l|r|}
\hline
\textbf{Componente} & \textbf{Disco} \\
\hline
Imagen del reconocimiento (AMD) & 83,5~GB \\
\hline
Caché de núcleos de cálculo compilados & 11,6~GB \\
\hline
Caché del modelo & 3,1~GB \\
\hline
Imagen de la aplicación & 0,35~GB \\
\hline
\textbf{Total aproximado} & \textbf{$\sim$100~GB} \\
\hline
\end{tabular}
\caption{Ocupación en disco del despliegue en contenedores sobre AMD, medida con
\texttt{docker system df -v} en el equipo de desarrollo.}
\label{tab:disco}
\end{table}

Lo llamativo es de dónde viene: \textbf{no es el modelo}, que son 2,85~GB, sino la imagen
base de la plataforma de cálculo de AMD. Inspeccionarla por dentro explica el motivo, y el
motivo es que el contenedor transporta código para tarjetas que el equipo no tiene.

Esa plataforma compila por anticipado los núcleos de cálculo para cada arquitectura de
tarjeta, y la imagen empleada los incluye \emph{todos}: su etiqueta declara literalmente
\texttt{device-all}. Medido dentro de ella, una sola de sus bibliotecas de álgebra contiene
núcleos para \textbf{veinticuatro arquitecturas distintas}, desde las de centro de datos
hasta tres generaciones de tarjetas de consumo; de sus 792~MB, los que corresponden a la
tarjeta empleada en este trabajo son \textbf{10~MB}. Otro componente aparece como trece
bibliotecas separadas, una por arquitectura, que suman 4,4~GB y de las que se usa una.

A ello se añaden dos circunstancias más. La imagen incluye el kit de \emph{desarrollo} de la
plataforma, 15~GB de compiladores y cabeceras que no hacen falta para ejecutar inferencia
pese a que la imagen se anuncia como de ejecución. Y los 11,6~GB de núcleos compilados en el
primer arranque se generan porque, para la generación de tarjeta concreta empleada, esa
plataforma no distribuye binarios precompilados y hay que producirlos en destino.

Conviene precisar el grado de confianza de las cifras. Las de la tabla proceden de
\texttt{docker system df -v}, que es lo que el sistema informa; el contenido real del sistema
de ficheros de la imagen, medido desde dentro, son 34~GiB. La diferencia se debe a cómo el
motor de contenedores contabiliza capas almacenadas y desplegadas, y no se ha podido acotar
con precisión por requerir privilegios de administrador sobre su directorio de datos. Para
planificar un despliegue conviene por tanto tomar la cifra alta, que es la conservadora.

Tres consecuencias prácticas. La primera, que un equipo con un disco pequeño, que es lo
habitual en un portátil de centro, puede quedarse sin espacio por esta vía y no por el
modelo. La segunda, que la ejecución nativa ocupa alrededor de 5~GB y evita el problema por
completo, lo que la convierte en una alternativa razonable y no en un apaño. La tercera, que
esto es reducible: partir de una imagen base restringida a la familia de tarjeta del equipo,
en lugar de a todas, eliminaría por construcción la mayor parte de lo que se transporta. No
se ha hecho, de modo que no se aporta la cifra resultante; queda como la mejora pendiente
más rentable del despliegue.

Hay una última observación que afecta menos al espacio y más a la reproducibilidad, y por eso
conviene dejarla escrita: la imagen base empleada es una compilación \textbf{nocturna} y no
una versión estable, según declaran sus propios parámetros de construcción. Para el
desarrollo es aceptable, porque era la vía de tener soporte de una tarjeta muy reciente; para
un despliegue en un centro no lo es, porque la referencia se mueve bajo los pies y dos
construcciones separadas en el tiempo no producen el mismo sistema. Fijar una versión estable
concreta es requisito previo a cualquier distribución, y es coherente con el criterio que
gobierna el resto del trabajo: un resultado que no se puede reconstruir no es un resultado.

Del perfil de NVIDIA no se da cifra porque no se ha medido: no se dispuso de una tarjeta de
ese fabricante. Cabe esperar que sea sustancialmente menor, porque su imagen base es de
ejecución y no de desarrollo, pero es una expectativa y no un dato.

\section{Cómo se distribuye el sistema}
\label{sec:contenedores}

\subsection*{Qué es un contenedor, y por qué se usa aquí}

El sistema se entrega \textbf{contenerizado}. Un contenedor es un paquete que incluye el
programa junto con todas sus dependencias, es decir el intérprete, las bibliotecas y sus
versiones exactas, de modo que se ejecuta igual en cualquier equipo que tenga el motor de
contenedores instalado, sin instalar nada más en el sistema anfitrión. Docker es la implementación
empleada aquí, y es la habitual.

La razón de recurrir a ello no es seguir una moda de despliegue: es que el entorno de este
sistema es especialmente incómodo de montar a mano. Requiere una versión concreta de la
biblioteca de cálculo, compilada contra la plataforma del fabricante de la tarjeta gráfica,
más el entorno de ejecución de la aplicación. Reproducir eso en el equipo de un centro sería
una tarde de trabajo para alguien con conocimientos, y una fuente permanente de averías al
actualizar el sistema operativo. Con contenedores, el centro instala Docker y ejecuta un
comando.

El sistema se reparte en \textbf{dos contenedores} y no en uno, porque tienen necesidades
distintas: el del reconocimiento necesita acceso a la tarjeta gráfica y pesa varios
gigabytes; el de la aplicación web no necesita tarjeta y es pequeño. Separarlos permite
además sustituir el modelo sin tocar la aplicación, que es la frontera descrita en la
sección~\ref{sec:integracion}.

\subsection*{Sobre qué sistemas operativos se puede ejecutar}

\begin{table}[H]
\centering
\footnotesize
\setlength{\tabcolsep}{4pt}
\begin{tabular}{|p{3.1cm}|c|c|p{5.6cm}|}
\hline
\textbf{Sistema} & \textbf{Vía} & \textbf{Acelerador} & \textbf{Estado} \\
\hline
Linux, tarjeta AMD & Contenedores & ROCm &
    \texttt{make docker-amd}. Es la plataforma de desarrollo del trabajo. \\
\hline
Linux, tarjeta NVIDIA & Contenedores & CUDA &
    \texttt{make docker-nvidia}. Requiere el conjunto de herramientas de contenedores del
    fabricante. \\
\hline
macOS, Apple Silicon & \textbf{Nativa} & MPS &
    \texttt{make demo}. En contenedor no alcanza la GPU; se explica más abajo. \\
\hline
Windows, como cliente & \multicolumn{3}{p{8.4cm}|}{\textbf{Funciona sin nada especial.} Es
    solo un navegador. Es además el caso más frecuente: los dispositivos del alumnado.} \\
\hline
Windows, como servidor & \multicolumn{3}{p{8.4cm}|}{No probado. Existe una vía a través del
    subsistema Linux de Windows, pero el descriptor de contenedores requiere cambios; se
    detalla más abajo.} \\
\hline
Sin tarjeta gráfica & \multicolumn{3}{p{8.4cm}|}{Arranca y transcribe, pero no sostiene el
    directo. Sirve para verificar la instalación y la red del aula, no para dar clase.} \\
\hline
\end{tabular}
\caption{Sistemas operativos sobre los que puede ejecutarse el servidor.}
\label{tab:vias}
\end{table}

\subsection*{De dónde salen las imágenes}

Conviene ser preciso en esto, porque determina lo que un centro tiene que hacer y lo que
todavía no puede hacer.

\textbf{Las imágenes de este sistema no están publicadas en ningún registro público.} No
existe, por tanto, una orden \texttt{docker pull} que las descargue: se construyen en el
propio equipo a partir del repositorio del proyecto, y eso es lo que hace la opción
\texttt{-{}-build} de los comandos de la tabla~\ref{tab:vias}. El proceso, la primera vez,
consiste en descargar las imágenes base públicas de las que se parte, instalar sobre ellas
las dependencias declaradas y copiar el código del proyecto.

Las imágenes base sí proceden de registros públicos, y son las que recoge la
tabla~\ref{tab:imagenes-base}:

\begin{table}[H]
\centering
\footnotesize
\begin{tabular}{|l|l|}
\hline
\textbf{Componente} & \textbf{Imagen base} \\
\hline
Reconocimiento, AMD & \texttt{rocm/pytorch} \\
\hline
Reconocimiento, NVIDIA & \texttt{pytorch/pytorch:2.5.1-cuda12.4-cudnn9-runtime} \\
\hline
Aplicación, compilación & \texttt{mcr.microsoft.com/dotnet/sdk:10.0} \\
\hline
Aplicación, ejecución & \texttt{mcr.microsoft.com/dotnet/aspnet:10.0} \\
\hline
\end{tabular}
\caption{Imágenes base de las que parten los contenedores del sistema.}
\label{tab:imagenes-base}
\end{table}

La versión concreta de la imagen base del reconocimiento se fija desde fuera y no dentro del
fichero de construcción, a propósito: cada combinación de tarjeta y controlador exige una
distinta, y codificarla obligaría a editar el fichero en cada despliegue.

\textbf{Qué implica que no estén publicadas.} Hoy, el centro necesita el repositorio del
proyecto y una máquina capaz de construir, lo que incluye descargar las dependencias de
compilación de la aplicación. Publicar las imágenes en un registro, por ejemplo el del propio
alojamiento del código, reduciría el despliegue a descargar el descriptor y ejecutar
\texttt{docker compose up}, sin repositorio ni construcción, y haría el primer arranque
mucho más rápido y mucho menos sensible al equipo. Es una mejora de distribución evidente y
queda como tarea pendiente, no como decisión: no se ha hecho porque el trabajo no ha llegado
a la fase de distribuir el sistema a terceros.

\subsection*{El caso de Windows}

Merece tratamiento propio por una razón práctica: es, con diferencia, el sistema más
extendido en los centros educativos españoles, de modo que despacharlo en una línea sería
ignorar el escenario más probable. Conviene separar dos papeles que se confunden con
facilidad.

\textbf{Como cliente no hay nada que hacer.} Los equipos del alumnado, y cualquier pantalla
del aula que solo muestre subtítulos, funcionan con Windows exactamente igual que con
cualquier otro sistema: abren una dirección en el navegador. Este es el caso más frecuente
en la práctica, y conviene decirlo explícitamente porque la advertencia siguiente puede
leerse mal.

\textbf{Como servidor no se ha probado, pero la vía existe.} El sistema no se ha desplegado
sobre Windows en este trabajo, y por tanto no se afirma que funcione. Sí puede describirse
con precisión qué haría falta, porque el obstáculo está localizado.

El camino sería el subsistema Linux de Windows, que hoy permite acceso a la tarjeta gráfica
desde dentro. Tanto la plataforma de NVIDIA \cite{microsoft2024cudawsl} como la de AMD
\cite{amd2025rocmwsl} exponen la tarjeta a través de un dispositivo propio del subsistema,
\texttt{/dev/dxg}, en lugar de los dispositivos que se usan en un Linux nativo. Y ahí está
justamente el problema: \textbf{el perfil de AMD de este proyecto pasa al contenedor
\texttt{/dev/kfd} y \texttt{/dev/dri}, que bajo el subsistema no existen}. Ese perfil
requeriría por tanto modificarse. El de NVIDIA está mejor situado, porque no enumera
dispositivos a mano sino que delega en el conjunto de herramientas del fabricante, que es
precisamente el que se encarga de exponer el dispositivo correcto dentro del contenedor.

Hay un segundo obstáculo, menor pero real: los contenedores comparten la red del anfitrión,
decisión tomada por un problema concreto de la distribución de desarrollo
(sección~\ref{sec:contenedores}). Ese modo se comporta de forma distinta en el entorno de
Docker para Windows, así que habría que volver a publicar puertos, que es lo que se hace
normalmente.

Queda una tercera vía que conviene no perder de vista, porque puede ser la más simple:
\textbf{ejecutar sin contenedores}. Los dos componentes del sistema son de por sí
multiplataforma: la biblioteca de cálculo del servicio de reconocimiento tiene versión para
Windows con tarjetas NVIDIA, y el entorno de ejecución de la aplicación es de Microsoft y
nativo allí. Lo que no es multiplataforma es la automatización del proyecto, que asume un
intérprete de órdenes de tipo Unix; habría que ejecutar los comandos a mano. Tampoco se ha
probado.

En resumen, y para que un centro sepa a qué atenerse: sobre Windows con tarjeta NVIDIA el
despliegue es plausible y probablemente cercano, sobre Windows con tarjeta AMD exige
modificar el descriptor, y en ambos casos habría que verificarlo antes de comprometerse.
Queda recogido como tarea pendiente en lugar de como afirmación.

\subsection*{El caso de Apple Silicon}

Merece explicación aparte porque es contraintuitivo y porque un centro con equipamiento
Apple lo encontrará de inmediato. Afecta a \textbf{toda la gama con chips de la serie M},
tanto portátiles como equipos de sobremesa, porque lo que lo provoca es la arquitectura del
procesador y no el formato del equipo.

Conviene señalarlo porque los equipos de sobremesa de esa gama son, de hecho, candidatos
razonables al papel de servidor de aula: son silenciosos, consumen poco y están pensados
para permanecer encendidos, que es exactamente lo que este sistema necesita. Un centro que
ya disponga de uno tiene una opción sensata, siempre que lo despliegue por la vía correcta.

Y la vía correcta no es la que parecería. La aplicación funciona sin cambios en esa
arquitectura, dentro o fuera de contenedor; el problema está en el servicio de
reconocimiento.

Ejecutado de forma nativa, el servicio detecta el acelerador integrado en el chip y lo
emplea, con una particularidad que el código resuelve solo: fuerza precisión simple, porque
Metal no implementa la precisión reducida para todas las operaciones del modelo y con ella
la inferencia falla o devuelve silencio.

Ejecutado en contenedor, en cambio, \textbf{queda en procesador y no hay forma de evitarlo}.
La causa no es del sistema desarrollado: en macOS los contenedores corren dentro de una
máquina virtual Linux que no expone el acelerador de Apple. Y quedarse en procesador
significa 1,6 veces la velocidad del audio (sección~\ref{sec:hardware}), insuficiente para
directo. De ahí que en estos equipos la vía correcta sea la nativa, aunque sea la menos
cómoda.

Queda una cosa por declarar, porque afecta a cuánta confianza merece esta recomendación:
\textbf{el rendimiento de ese acelerador no se ha medido en este trabajo}. Se ha verificado
que el servicio lo detecta y transcribe correctamente, pero el barrido de dispositivos
comparó únicamente tarjeta gráfica dedicada y procesador (sección~\ref{sec:hardware}).
Además, dentro de la propia gama el rendimiento varía mucho: no es comparable el chip de
entrada de un portátil con las variantes de gama alta de un sobremesa, que multiplican el
número de núcleos gráficos. Un centro que vaya a desplegar sobre este equipamiento debería
medir el factor de tiempo real en su equipo concreto, con el mismo procedimiento del
experimento, antes de comprometerse.

\subsection*{Estado de verificación de cada vía}

Por el mismo motivo se declara aquí, y no en una nota al pie, qué está probado y qué no:

\begin{itemize}
  \item \textbf{Vía nativa sobre Linux con tarjeta AMD:} es la empleada durante todo el
        desarrollo y la que produjo las mediciones de esta memoria. Probada.
  \item \textbf{Vía en contenedores sobre Linux con tarjeta AMD:} \textbf{construida y
        ejecutada}. Los dos contenedores arrancan, superan su comprobación de salud y se
        mantienen en servicio; las cifras de ocupación en disco de la
        tabla~\ref{tab:disco} proceden de ese despliegue. Lo que no se ha hecho es una
        sesión de aula completa a través de ellos.
  \item \textbf{Vía en contenedores sobre Linux con tarjeta NVIDIA:} sin probar, por no
        disponer de una tarjeta de ese fabricante. El descriptor está escrito y difiere del
        anterior solo en la imagen base y en el mecanismo de acceso a la tarjeta.
  \item \textbf{Vía nativa sobre Apple Silicon:} el servicio detecta el acelerador y
        transcribe. Sin medición de rendimiento.
  \item \textbf{Windows, en cualquier variante:} sin probar como servidor. Como cliente
        funciona igual que cualquier otro sistema.
\end{itemize}

Este inventario es deliberado, y su nivel de detalle también. El
capítulo~\ref{cap:desarrollo} documenta varios casos en los que algo funcionaba en
apariencia y producía resultados falsos; dar por verificada una vía de despliegue que nadie
ha ejecutado sería exactamente el mismo error, y darla por no verificada cuando sí se ha
ejecutado tampoco ayuda a quien tenga que desplegarla.

\section{Puesta en marcha}

\subsection*{Paso 1: obtener el proyecto}

Basta con el repositorio; no hay que descargar el modelo a mano, porque el contenedor lo
hace en el primer arranque y lo conserva.

\subsection*{Paso 2: configurar la clave del puesto docente}

Se define en un fichero \texttt{.env} junto al descriptor de los contenedores:

\begin{quote}\ttfamily
CLAVE\_DOCENTE=la-clave-que-elija-el-centro
\end{quote}

\textbf{Este paso no es opcional, y el descriptor de contenedores lo impone.} Sin clave, la
vista de emisión queda abierta y cualquier alumno conectado a la red podría tomarla y
emitir por el aula. Durante un tiempo la variable tenía el valor vacío por defecto, de modo
que seguir estos pasos sin definirla levantaba el sistema desprotegido sin ninguna
advertencia; ahora el arranque se detiene y explica qué falta. La interfaz sigue avisando
además cuando la aplicación corre sin clave, que es el modo de desarrollo.

En el mismo fichero conviene fijar la política de reinicio:

\begin{quote}\ttfamily
REINICIO=unless-stopped
\end{quote}

Con ella, los contenedores vuelven a levantarse solos cuando arranca el equipo, que es lo
que se quiere de un servidor de aula: tras un corte de luz o un reinicio nocturno, el
sistema está disponible sin que nadie tenga que entrar a arrancarlo. El valor por defecto es
el contrario, \texttt{no}, porque en un equipo de desarrollo unos contenedores que resucitan
en cada encendido son una molestia, y el del reconocimiento además mantiene ocupado un
núcleo y varios gigabytes de memoria mientras vive.

Se emplea \texttt{unless-stopped} y no \texttt{always} por una diferencia que importa en un
centro: si alguien detiene el sistema a propósito, debe seguir detenido hasta que se decida
arrancarlo, y no reaparecer en el siguiente encendido.

\subsection*{Paso 3: levantar el sistema}

Un único comando, elegido según la tabla~\ref{tab:vias}:

\begin{quote}\ttfamily
make docker-amd \qquad\qquad\rmfamily(Linux, tarjetas AMD, plataforma ROCm)\\
\ttfamily make docker-nvidia \qquad\rmfamily(Linux, tarjetas NVIDIA, plataforma CUDA)\\
\ttfamily make demo \qquad\qquad\quad\rmfamily(macOS con Apple Silicon, ejecución nativa)
\end{quote}

Entre las dos primeras la diferencia es únicamente la imagen base del servicio de
reconocimiento; el código es idéntico, porque la biblioteca de cálculo expone la misma
interfaz sobre las dos plataformas. La tercera no usa contenedores por el motivo explicado
en la sección~\ref{sec:contenedores}, y requiere tener preparado el entorno de Python y
el de la aplicación en el propio equipo.

\textbf{El primer arranque es lento y eso es normal.} Hay que descargar la imagen base, que
en el caso de AMD son más de ochenta gigabytes (tabla~\ref{tab:disco}), descargar el modelo
y compilar los núcleos de cálculo de la tarjeta. Todo ello queda guardado en volúmenes, de
modo que los arranques siguientes son inmediatos. Conviene por tanto hacer el primer
arranque con un día de margen sobre la primera clase, y con una conexión que aguante esa
descarga.

\subsection*{Paso 4: comprobar que funciona}

El comando anterior imprime al arrancar las dos direcciones del sistema, recogidas en la
tabla~\ref{tab:direcciones}. Son distintas y no son intercambiables:

\begin{table}[H]
\centering
\footnotesize
\begin{tabular}{|l|l|p{6.2cm}|}
\hline
\textbf{Quién} & \textbf{Dirección} & \textbf{Observación} \\
\hline
Docente & \texttt{http://localhost:5203/broadcast} &
    Debe abrirse \textbf{en el propio equipo servidor}. \\
\hline
Alumnado & \texttt{http://<ip-del-equipo>:5203/view} &
    Desde cualquier dispositivo de la misma red. \\
\hline
\end{tabular}
\caption{Direcciones de acceso al sistema.}
\label{tab:direcciones}
\end{table}

La restricción de la primera fila tiene una causa concreta y conviene entenderla antes de
intentar sortearla: los navegadores solo conceden acceso al micrófono en conexiones seguras
o en la máquina local. Como el sistema opera dentro de la red del centro y sin certificado,
\textbf{la emisión tiene que hacerse desde el equipo servidor}. Si se quisiera emitir desde
otro equipo del aula habría que servir la aplicación por HTTPS con un certificado
reconocido por los dispositivos del centro, lo que es perfectamente posible pero excede lo
que este trabajo despliega.

\subsection*{Paso 5: dar clase}

En la vista del docente: introducir la clave, pegar en el recuadro del glosario los términos
de la sesión, uno por línea, y pulsar el botón de comenzar. Los subtítulos aparecen a la vez
en la pantalla del docente y en la de todo el alumnado conectado.

El glosario conviene tomarlo del temario o de las transparencias de la sesión. No hace falta
que sea exhaustivo: sirve para que los términos propios de la asignatura se resalten
visualmente en el subtítulo.

\subsection*{Paso 6: detener el sistema}

\begin{quote}\ttfamily
make docker-parar
\end{quote}

Detiene los contenedores sin borrar el modelo descargado ni las cachés, de modo que el
arranque siguiente vuelva a ser inmediato.

\section{Incorporar el modelo adaptado}

El resultado del ajuste fino se activa añadiendo una línea al mismo fichero \texttt{.env}:

\begin{quote}\ttfamily
ADAPTADOR=/adaptador
\end{quote}

No hay que reconstruir nada ni modificar la aplicación: el servicio de reconocimiento aplica
el adaptador sobre el modelo base al cargarlo. Dejar la variable vacía sirve el modelo sin
adaptar. Es la comprobación práctica de la frontera descrita en la
sección~\ref{sec:integracion}.

\section{Qué hacer si algo no funciona}

Los tres fallos que se han observado en la práctica tienen la misma característica que los
incidentes del capítulo~\ref{cap:desarrollo}: el sistema no se cae, simplemente funciona
peor sin decirlo.

\textbf{Los subtítulos tardan mucho y llegan en bloques largos.} Mirar el contador de la
vista del docente que separa los cortes cerrados \emph{por pausa} de los cerrados \emph{por
agotar el tope}. Si dominan los segundos, el detector de silencios no está encontrando las
pausas. La causa habitual no está en el sistema sino en el navegador: el control automático
de ganancia amplifica el ruido de fondo cuando el hablante calla, y la energía nunca baja lo
suficiente. La aplicación lo desactiva al capturar; si aun así ocurre, revisar que no haya
un procesado de audio activo en el sistema operativo o en el propio micrófono.

\textbf{Aparecen frases que nadie ha dicho.} El sistema filtra los dos patrones
reconocibles, las muletillas de subtítulos de vídeo y los bucles de repetición, pero no
puede filtrar una alucinación verosímil (sección~\ref{sec:filtro}). Si son frecuentes, casi
siempre indica que al micrófono le llega poca voz y mucho fondo: acercarlo al docente antes
que tocar ninguna configuración.

\textbf{La calidad cae de golpe a mitad de clase.} La degradación frente al ruido no es
gradual: el sistema pasa de funcionar bien a fallar en un margen de cinco decibelios
(sección~\ref{sec:ruido}). Si el aula se ha animado, el remedio es el micrófono de
proximidad, no un ajuste del software.

\textbf{El contenedor de reconocimiento no arranca con tarjeta AMD.} Casi siempre son los
identificadores de los grupos \texttt{video} y \texttt{render}, que varían entre equipos y
están fijados numéricamente en el descriptor. Se consultan con
\texttt{getent group video render} y se ajustan.

\section{Arranque sin contenedores}

Los dos atajos que se emplearon durante el desarrollo sirven también para desplegar, y son
la vía obligada en macOS:

\begin{quote}\ttfamily
make app \qquad\rmfamily arranca solo la aplicación, con un motor de reconocimiento
simulado\\
\ttfamily make demo \qquad\rmfamily arranca el servicio de reconocimiento real y la
aplicación, y limpia al salir\\
\ttfamily make demo-lora \quad\rmfamily igual, pero sirviendo el modelo adaptado
\end{quote}

El primero merece atención porque resuelve un problema práctico de despliegue: permite
comprobar que la red del aula, los dispositivos del alumnado y las direcciones de acceso
funcionan \emph{antes} de que entre en juego el reconocimiento. Si los subtítulos simulados
llegan a los móviles de la clase y los reales no, el problema está en el servicio de
reconocimiento; si no llegan ni los simulados, está en la red. Separar esas dos causas
ahorra la mayor parte del tiempo de diagnóstico.
